\documentstyle[%


righttag%


,11pt%


,amssymb%


,russian%               remove in Eng.


,izvru%                 Set izven% in Eng.


]{amsart}





%       Math definitions


\newcommand{\g}{\Gamma}


\newcommand{\Pd}[2]{\frac{\partial{#1}}{\partial{#2}}}


\newcommand{\pdd}[1]{\frac{\partial}{\partial{#1}}}


\newcommand{\grd}{\text{grad}\,\rho}


\newcommand{\ovl}{\overline}


\newcommand{\zt}{\zeta}


\newcommand{\ztz}{(\zt,z)}


\newcommand{\z}{(z)}


\newcommand{\ig}{\int_\Gamma}


\newcommand{\igs}{\int_{\Gamma\cap S}}


\newcommand{\ld}{\land}


\newcommand{\sm}{\sum\limits_{k=1}^n}


\newcommand{\r}[1]{\rho_{#1}}


\newcommand{\ovlr}[1]{\rho_{\ovl{#1}}}


\newcommand{\bm}{Бох\-не\-ра--Мар\-ти\-нел\-ли }


\newcommand{\tts}[1]{}


%\renewcommand{\baselinestretch}{0.96}





%\hoffset=-15mm%                  For Eng.


 


\setcounter{page}{37}


\god{1999}


\nomer{6~(445)} %                        Remove in Eng.


%\nomer{Vol.\,43, No.\,6}%               Set in Eng.


%\str{\pageref{begin}--\pageref{end}}%  Set in Eng.


\udk{517.55}





\author{\it А.М.\,Кытманов, О.В.\,Ходос}


% NB:   ^^ remove in Eng.


\title{Об условиях голоморфного продолжения гладких\\ $CR$-функций


в фиксированную область}


\thanks{Работа выполнена при финансовой поддержке


Российского фонда фундаментальных исследований (грант 96-01-00080).}





\date{}


%\pagestyle{myheadings}%         Set in Eng.


\pagestyle{plain} %              Remove in Eng.


\begin{document}


%\thispagestyle{plain}%          Set in Eng.


\maketitle


\label{begin}





Пусть $\Omega$ --- ограниченная область в $\Bbb C^n$ с гладкой


(класса $C^1$) границей $\partial\Omega$, а $\Gamma$ --- гладкая 


ориентируемая гиперповерхность в окрестности $U$ замыкания $\ovl\Omega$,


пересекающая $S=\partial\Omega$ трансверсально, вида


$$


\Gamma=\{z\in\ovl\Omega:\rho(z)=0\},\quad \rho\in C^1(U),


\quad\grd\ne0 \text{\; на \;}\Gamma.


$$


Пусть $\Gamma$ делит $\Omega$ на


области $\Omega^+=\{z:\rho(z)>0\}$ и $\Omega^-=\{z:\rho(z)<0\}$. Ориентация


$\Gamma$ согласована с $\Omega^+$.





В [1] и [2] рассматривалась задача о нахождении условий на


$CR$-функцию $f$ (определенную на $\Gamma$), при которых $f$ голоморфно 


продолжается в $\Omega^+$. Эти условия связаны с гармоническим продолжением 


интеграла \bm{} от функции $f$ или с вещественно-аналитическим продолжением


интеграла Коши--Фантаппье. Приведем один из результатов, который прямо 


следует из теорем работ [1], [2].





Пусть $\Omega$ --- единичный шар с центром в нуле, $0\in\Omega^-$. Ядром


\bm{} является


$$


U\ztz=\sm(-1)^{k-1}\frac{\partial g}{\partial\zeta_k}\,d\ovl\zeta[k]\land


d\zeta,


$$


где $d\zeta=d\zeta_1\land\dotsb\land d\zeta_n$, $d\ovl\zeta[k]=d\ovl\zeta_1


\land\dotsb\land d\ovl\zeta_{k-1}\land d\ovl\zeta_{k+1}\land\dotsb\land


d\ovl\zeta_n$,


$$


g\ztz=


\begin{cases}


-\dfrac{(n-2)!}{(2\pi i)^n}|\zt-z|^{2-2n}&\text{для $n>1$};\\


\phantom{-}\dfrac1{2\pi i}\ln|\zt-z|^2&\text{для $n=1$}


\end{cases}


$$


--- фундаментальное решение уравнения Лапласа.





\begin{thm}


Пусть $CR$-функция $f\in\cal L^1 (\Gamma)$ и


$$


c_\alpha=\bigg[\frac{\partial^\alpha}{\alpha!}\int_\Gamma


f(\zeta)U(\zeta,z)\bigg]\bigg|_{z=0},


$$


где $\alpha=(\alpha_1,\dotsc,\alpha_n)$ --- мультииндекс, $\alpha!=\alpha_1!


\dotsb\alpha_n!$, $\partial^\alpha=\frac{\partial^{\|\alpha\|}}


{\partial z_1^{\alpha_1}\dotsb\partial z_n^{\alpha_n}}$, $\|\alpha\|=


\alpha_1+\dotsb+\alpha_n$. Для того чтобы $f$ голоморфно продолжалась в


$\Omega^+$, необходимо и достаточно, чтобы


\begin{equation}


\underset{\|\alpha\|\to\infty}{\varlimsup}{}


\sqrt[\uproot{3}\|\alpha\|]{|c_\alpha|d_\alpha}


\leqslant1,


\end{equation}


где $d_\alpha=\max\limits_{\ovl\Omega}|z^\alpha|$, $z^\alpha=


z_1^{\alpha_1}\dotsb z_n^{\alpha_n}$.


\end{thm}


\noindent Заметим, что при $n=1$ всякая функция $f$ является $CR$-функцией.





Нашей целью является получение условий (1) не в терминах роста производных


интеграла \bm, а в терминах роста производных самой функции $f$.


Для этого, используя формулу Стокса, преобразуем выражение для коэффициентов


$c_{\alpha}$, перебрасывая производные с ядра Бохнера--Мар\-ти\-нел\-ли на


саму функцию $f$. При этом получаются дополнительно интегралы по множеству


$\Gamma\cap S$.


Тем самым условие (1) для бесконечно дифференцируемых функций будет более


конструктивным, чем в работах [1], [2], т.\,к. даст возможность проверять


его для различных классов конкретных гладких функций с известным поведением


производных на поверхности $\Gamma$.





Сначала  приведем формулы для производных интеграла \bm, обобщающие


соответствующие формулы из ([3], \S\,4, c.\,34).





Пусть


$$


(Mf)(z)=\int_\Gamma f(\zeta)U(\zeta,z),\quad z\in\Omega^\pm.


$$





\begin{lem}


Если $f\in C^1(\Gamma)$, то


\begin{multline*}


\frac{\partial\,Mf}{\partial z_m}


=\ig\Pd{f}{\zt_m} U+(-1)^{n+m}\sm\ig\Pd{f}{\ovl\zt_k}


\Pd{g}{\zt_k}\,d\ovl\zt\ld d\zt\,[m]+\\


+(-1)^{n+m}\sm(-1)^k\igs f\Pd{g}{\zt_k}\,d\ovl\zt\,[k]\ld d\zt[m],


\quad z\in\Omega\backslash\Gamma.


\end{multline*}


\end{lem}





Обозначим


$\r{k}=\displaystyle\Pd{\rho}{\zt_k}\frac1{|\grd|}$, $\ovlr{k}=\ovl\rho_k$.





\begin{lem}


Пусть для $z\notin\g$


$$


\Phi\z=i^n2^{n-1}\ig f(\zt)g\ztz\,d\sigma(\zt)


$$


--- потенциал простого слоя, тогда справедлива формула


\begin{multline*}


\Pd{\Phi}{z_m}=-\ig f\r{m}U\ztz


+i^n2^{n-1}\sm\ig\bigg(\r{k}\pdd{\zt_m}(f\ovlr{k})-\r{m}\pdd{\zt_k}


(f\ovlr{k})\bigg)g\ztz\,d\sigma(\zt)+ \\


+\igs\bigg(\sum_{k=1}^{m-1}(-1)^{n+m+k}\ovlr{k}\,d\zt[k,m]\ld d\ovl\zt


+\sum_{k=m+1}^n(-1)^{n+m+k-1}\ovlr{k}\,d\zt[m,k]\ld


d\ovl\zt\bigg)f(\zt)g\ztz.


\end{multline*}


\end{lem}





{\bf Предложение.}


{\it Если $\g\in C^2$ и $f\in C^2(\Gamma)$, то}


\begin{multline}


\Pd{\,Mf}{z_m}


=\ig\bigg(\Pd{f}{\zt_m}-\r{m}\sm\r{k}\Pd{f}{\ovl\zt_k}\bigg)U(\zeta,z)+


2^{n-1}i^n\ig g(\zeta,z)A_nd\sigma+\\


+\igs g(\zeta,z)\r{m}\sm\Pd{f}{\ovl\zt_k} B_n


+\igs\sm(-1)^{n+m+k} f\Pd{g}{\zt_k}\,d\ovl\zt[k]\ld d\zt[m],


\end{multline}


{\it где}


\begin{align*}


A_n&=\sum_{k,l=1}^n\bigg(\r{l}\pdd{\zt_k}\Big(\Pd{f}{\ovl\zt_k}


\r{m}\ovlr{l}\Big)


-\r{k}\pdd{\zt_l}\Big(\Pd{f}{\ovl\zt_k}\r{m}\ovlr{l}\Big)\bigg),\\


%


B_n&=\sum_{l=1}^{k-1}(-1)^{k+l}\ovlr{l}  \,d\ovl\zt\ld d\zt[l,k]


+\sum_{l=k+1}^n(-1)^{k+l-1}\ovlr{l}\,d\ovl\zt\ld d\zt[k,l].


\end{align*}





Пусть $f$ --- $CR$-функция, т.\,е. $\ovlr{l}\Pd{f}{\ovl\zt_k}=\ovlr{k}


\Pd{f}{\ovl\zt_l}$ для всех $k,l=1,\dotsc,n$. Тогда в (2) второе и


третье слагаемые равны нулю, т.\,к.


\begin{align*}


A_n&=\sum_{k,l=1}^n\r{l}\pdd{\zt_k}\bigg(\Pd{f}{\ovl\zt_k}\r{m}\ovlr{l}\bigg)


-\sum_{k,l=1}^n\r{l}\pdd{\zt_k}\bigg(\Pd{f}{\ovl\zt_l}\r{m}\ovlr{k}\bigg)=\\


&=\sum_{k,l=1}^n\r{l}\pdd{\zt_k}\bigg(\r{m}\Big(\Pd{f}{\ovl\zt_k}\ovlr{l}


-\Pd{f}{\ovl\zt_l}\ovlr{k}\Big)\bigg)=0,\\


%


B_n&=\sm\sum_{l=1}^{k-1}(-1)^{k+l}


\Pd{f}{\ovl\zt_k}\ovlr{l}\,d\ovl\zt\ld d\zt[l,k]


+\sm\sum_{l=1}^{k-1}(-1)^{k+l-1}


\Pd{f}{\ovl\zt_l}\ovlr{k}\,d\ovl\zt\ld d\zt[l,k]=\\


%


&=\sm\sum_{l=1}^{k-1}(-1)^{k+l}\bigg(\Pd{f}{\ovl\zt_k}\ovlr{l}


-\Pd{f}{\ovl\zt_l}\ovlr{k}\bigg)d\ovl\zt\ld d\zt[l,k]=0.


\end{align*}


Таким образом, получаем





\begin{cor}


Пусть $\g\in C^2$, $f\in C^2(\Gamma)$, $f$ --- $CR$-функция. Тогда


$$


\Pd{\,Mf}{z_m}=\ig\bigg(\Pd{f}{\zt_m}


-\r{m}\sm\r{k}\Pd{f}{\ovl\zt_k}\bigg)U(\zeta,z)+


\igs\sm(-1)^{n+m+k} f\Pd{g}{\zt_k}\,d\ovl\zt[k]\ld d\zt[m],


\quad z\notin\Gamma.


$$


\end{cor}





Выберем функцию $\rho$ так, что $|\grd|\equiv1$  в окрестности $\Gamma$,


т.\,е. $\rho_k=\frac{\partial\rho}{\partial\zeta_k}$, а $\rho_{\ovl k}=


\frac{\partial\rho}{\partial\ovl\zeta_k}$. Это всегда можно сделать для


гиперповерхности $\Gamma$ класса $C^2$ ([4], \S\,2; [5]).


Введем оператор


\begin{equation}


D_m=\frac\partial{\partial\zeta_m}-\rho_m\sum_{k=1}^n\rho_k\frac\partial


{\partial\ovl\zeta_k}.


\end{equation}





\begin{lem}


Для $CR$-функций  $f$ справедливо тождество


$$


(D_m\circ D_l)(f)=(D_l\circ D_m)(f).


$$


\end{lem}





{\bf Доказательство.} Действительно,


\begin{multline*}


D_m\circ D_l


=\frac{\partial^2}{\partial\zeta_m\partial\zeta_l}


-\rho_l\sum_{k=1}^n\rho_k\frac{\partial^2}{\partial\ovl\zeta_k\partial\zeta_m}


-\frac{\partial}{\partial\zeta_l}\bigg(\rho_m\sum_{k=1}^n\rho_k\frac{\partial}


{\partial\ovl\zeta_k}\bigg)+


% previous partition ^


\rho_m\sum_{k=1}^n\rho_{kl}\frac{\partial}{\partial\ovl\zeta_k}


-\rho_l\sum_{k=1}^n\rho_{km}\frac{\partial}{\partial\ovl\zeta_k}+ \\


% newer partition


+\rho_l\sum_{k=1}^n\rho_{k}\frac{\partial}{\partial\ovl\zeta_k}


\bigg(\rho_m\sum_{s=1}^n\rho_l\frac{\partial}{\partial\ovl\zeta_s}\bigg)-


% removed partition


\rho_l\sum_{k=1}^n\rho_k\rho_{m\ovl k}\sum_{s=1}^n\rho_s\frac{\partial}


{\partial\ovl\zeta_s}+\rho_m\sum_{k=1}^n\rho_k\rho_{l\ovl k}


\sum_{s=1}^n\rho_s\frac\partial{\partial\ovl\zeta_s}=\\


=D_l\circ D_m+\rho_m\sum_{k=1}^n\bigg(\rho_{kl}


\frac{\partial}{\partial\ovl\zeta_k}


+\rho_k\rho_{l\ovl k}\sum_{s=1}^n\rho_s\frac{\partial}{\partial\ovl\zeta_s}


\bigg)-


\rho_l\sum_{k=1}^n\bigg(\rho_{km}\frac{\partial}{\partial\ovl\zeta_k}


+\rho_k\rho_{m\ovl k}\sum_{s=1}^n\rho_s\frac{\partial}{\partial\ovl\zeta_s}


\bigg);


\end{multline*}


здесь $\rho_{kl}=\frac{\partial^2\rho}{\partial\zeta_k\partial\zeta_l}$,


а $\rho_{k\ovl l}=\frac{\partial^2\rho}{\partial\zeta_k\partial\ovl\zeta_l}$.





Так как $\sum\limits_{k=1}^n\rho_k\rho_{\ovl k}\equiv1$ в некоторой окрестности


гиперповерхности $\Gamma$, следовательно,\linebreak


$\frac{\partial}{\partial z_s}\Big(\sum\limits_{k=1}^n\rho_k\rho_{\ovl k}


\Big)=0$


в этой окрестности, т.\,е.


$$


\sum_{k=1}^n\rho_{ks}\rho_{\ovl k}+\sum_{k=1}^n\rho_k\rho_{\ovl ks}=0,\quad


\sum_{k=1}^n\rho_k\rho_{\ovl ks}=-\sum_{k=1}^n\rho_{ks}\rho_{\ovl k}.


$$


Тогда


$$


D_m\circ D_l-D_l\circ D_m=


\rho_m\sum_{k=1}^n\bigg(\rho_{kl}\frac{\partial}{\partial\ovl\zeta_k}


-\rho_{kl}\rho_{\ovl k}\sum_{s=1}^n\rho_s\frac{\partial}{\partial\ovl\zeta_s}


\bigg)-


\rho_l\sum_{k=1}^n\bigg(\rho_{km}\frac{\partial}{\partial\ovl\zeta_k}


-\rho_{km}\rho_{\ovl k}\sum_{s=1}^n\rho_s\frac{\partial}{\partial\ovl\zeta_s}


\bigg).


$$


Так как $\rho_{\ovl k}\frac{\partial f}{\partial\ovl\zeta_s}=\rho_{\ovl s}


\frac{\partial f}{\partial\ovl\zeta_k}$ для $CR$-функций $f$, то


\begin{multline*}


D_m\circ D_l-D_l\circ D_m=


\rho_m\sum_{k=1}^n\bigg(\rho_{kl}\frac{\partial}{\partial\ovl\zeta_k}


-\rho_{kl}\sum_{s=1}^n\rho_{\ovl s}\rho_s\frac{\partial}{\partial\ovl\zeta_k}


\bigg)-


\rho_l\sum_{k=1}^n\bigg(\rho_{km}\frac{\partial}{\partial\ovl\zeta_k}


-\rho_{km}\sum_{s=1}^n\rho_{\ovl s}\rho_s\frac{\partial}{\partial\ovl\zeta_k}


\bigg)=\\


=\rho_m\sum_{k=1}^n\bigg(\rho_{kl}\frac{\partial}{\partial\ovl\zeta_k}


-\rho_{kl}\frac{\partial}{\partial\ovl\zeta_k}\bigg)-


\rho_l\sum_{k=1}^n\bigg(\rho_{km}\frac{\partial}{\partial\ovl\zeta_k}


-\rho_{km}\frac{\partial}{\partial\ovl\zeta_k}\bigg)=0.\quad\square


\end{multline*}





Перейдем к нахождению производной произвольного порядка от


интеграла Бохнера--Мар\-ти\-нел\-ли $Mf$. Пусть теперь $\Omega$ --- единичный


шар с центром в нуле.





\begin{thm}


Пусть $\g\in C^\infty$, $f\in C^\infty


(\Gamma)$, $f$ --- $CR$-функция. Тогда


\begin{multline}


\partial^\alpha Mf\big|_{z=0}=\int_\Gamma D^\alpha(f)\,


U(\zeta,0)+\frac{1}{(2\pi i)^n}\sum_{k=1}^n(-1)^{n+k-1}


\sum_{s=1}^n\sum_{l_s=1}^{\alpha_s}


(n+l_s+\alpha_{s+1}+\dotsb+\alpha_n-2)!\times\\


%


\times\int_{\Gamma\cap S}D_1^{\alpha_1}\circ\dotsb\circ


D^{\alpha_{s-1}}_{s-1}\circ D^{\alpha_s-l_s}_s\,(f)


\ovl\zeta^{\;l_s+\alpha_{s+1}+\dotsb+\alpha_n}\,d\,\ovl\zeta[k]\land


d\zeta[s].


\end{multline}


Здесь


$$


D^\alpha=D_1^{\alpha_1}\circ\dotsb\circ D^{\alpha_n}_n,\quad


D_j^{\alpha_j}=


\underbrace{D_j\circ\dotsb\circ D_j}_{\alpha_j \text{  раз}},


$$


$D_j$ задано в $(3)$, $j=1,\dotsc,n.$


\end{thm}





Таким образом, в условии (1) константу $c_\alpha$ можем брать из формулы (4).





Рассмотрим случай $n=1$. Введем обозначения


$$  D^m=\underbrace{D\circ\dots\circ D}_{m \text{ раз}},\quad


D=\frac{d}{d\zt}-\rho_1^2\frac{d}{d\ovl\zt}


$$


--- производная по касательной к $\Gamma$, $\rho_1=\dfrac{d\rho}{d\zt}$.





\begin{cor}


Если кривая $\g\in C^\infty(\Omega)$, $\xi$,


$\eta$ --- точки пересечения $\g$ с окружностью $S=\partial\Omega$,


$f\in C^\infty(\g)$, то


$$


\frac{d^m Mf}{dz^m}\Big|_{z=0}=\frac{1}{2\pi i}\ig \frac{D^m(f)}{\zeta}\,


d\zeta+


\frac{1}{2\pi i}\sum_{l=1}^m


(l-1)!(D^{m-l}(f)\big|_{\zeta=\xi}\ovl\xi^{\,l}


-D^{m-l}(f)\big|_{\zeta=\eta}\ovl\eta^{\,l}).


$$


\end{cor}





\begin{thebibliography}{9}





\bibitem{1}


Айзенберг Л.А., Кытманов А.М.


{\em О возможности голоморфного продолжения функций, заданных


на связном куске ее границы} //


Матем. сб. -- 1991. -- T.\,182. -- \No\,5. -- C.\,490--507.


\bibitem{2}


Айзенберг Л.А., Кытманов А.М.


{\em О возможности голоморфного продолжения функций, заданных


на связном куске ее границы}. II // Матем. сб. -- 1993. --


T.\,184. -- \No\,1. -- C.\,1--14.


\bibitem{3}


Кытманов А.М.


{\em Интеграл Бохнера--Мартинелли и его применения}. --


Новосибирск: Наука, 1992. -- 240 c.


\bibitem{4}


Волков Е.А.


{\em О границах подобластей, весовых классах Г\"ельдера и решениях


в этих классах уравнений Пуассона} //


Тр. Матем. ин-та АН СССР. -- 1972. -- Т.\,117. -- C.\,75--99.


\bibitem{6}


Горенский Н.Ю.


{\em Некоторые приложения дифференциальных свойств функции


расстояния до открытых множеств в $\Bbb C^n$}// 


Некотор. свойства голоморфн. функций


многих комплексн. переменных. -- Красноярск, 1973. -- C.\,203--208.





\end{thebibliography}





\vskip 0.5cm





\post{\it Красноярский государственный}{\it Поступила}


\post{\it технический университет}{14.08.1996}





\label{end}


\end{document}





